Many cryptographic algorithms (\Cref{sec:cryptography_basics:algorithms}) rely on one-way and trapdoor functions build upon mathematical structures (\Cref{sec:mathematical_underpinnings}): a \textbf{one-way} function is easy to compute but (computationally) hard to invert --- that is, it is prohibitively expensive to obtain the input given the corresponding output --- while a trapdoor function is easy to compute but (computationally) hard to invert without knowledge of additional information called \textbf{trapdoor} --- trapdoor functions are a specific instance of one-way functions. Roughly, hash functions (\Cref{sec:hash_functions}) rely on one-way functions, asymmetric cryptography (\Cref{sec:asymmetric_cryptography}) relies on trapdoor functions, and symmetric cryptography (\Cref{sec:symmetric_cryptography}) usually relies on (keyed pseudo)random permutations or (pseudo)random functions which do not fully match the definition of trapdoor functions. 

\exampleBlock{Real-world analogies for one-way and trapdoor functions}{A one-way function is like breaking an egg, shredding a sheet of paper, or mixing two colors of paint. Conversely, a trapdoor function is like a padlock, a sealed mailbox, or an hotel room door.}

At the time of writing (2025), no function was mathematically proven to be a one-way or trapdoor function. Still, there exist several candidates for trapdoor functions; we briefly describe the most relevant trapdoor functions below. Our goal is not to explain these trapdoor functions exhaustively, but rather to provide the reader with the basic knowledge necessary to connect these trapdoor functions to the algorithms described in later notes and to gain at least an elementary understanding of the subject.
